\documentclass[11pt]{article}
\usepackage{amsmath,amssymb,booktabs,geometry}
\geometry{margin=1in}

\title{Step 200: RH Iteration Arc 173--200 Closeout}
\author{six-birds-foundations-iii}
\date{}

\begin{document}
\maketitle

\section*{Verdict}

\[
\boxed{\texttt{V\_arc\_terminus\_declared}.}
\]

\section*{Theorem}

\textbf{Theorem (RH Track Iteration Arc 173--200 Terminus).}
Under the inherited Six Birds Foundations III framework discipline, inherited
records through Step 162, and the manager-led iteration arc Steps 173--200:

\begin{enumerate}
\item The retained typed-condition catalog for RH-analogous closure carriers
consists of the Cascade Reduction Theorem, CTMT, CRCFT, the Framework
RH-Carrier Dichotomy, and the Bridge Impossibility Corollary.

\item Every surveyed classically-RH-equivalent carrier exhibits CRCFT
foreclosure in one of the three modes CTMT, target-equivalence, or
bridge-failure.  No CRCFT coverage refuter was found.

\item Every surveyed non-CRE but RH-analogous carrier with a proved native
ledger yields native closure of its analog residual.

\item The Bridge Impossibility Corollary forecloses non-circular
non-CRE-to-Riemann transport under the framework discipline.  Path 2 is
structurally closed.

\item The only non-foreclosed Riemann attack path under the framework is Path
1: resolve a CRCFT-stuck record by external classical theorem.  In the current
arc this is the Burnol \(\kappa\) problem, equivalently the projected Sonine
kernel \(P_{L_a^\Gamma}K_a^{\Gamma,\mathrm{amb}}\).

\item Two substantial numerical experiments were produced: Branch C
nonvanishing of \(L_{\rho,k}(G)\) across 18 triples, and the
Beurling--Nyman harmonic-chain extension to \(N=2000\), where
\(\delta^2\log N\approx0.045\).

\item The cumulative framework deliverable is the audit document
\[
\texttt{anti\_loc/RH\_framework\_audit.md},
\]
with 925 lines, plus the step-local closeout artifacts.
\end{enumerate}

\textbf{Consequence.}
The arc reaches diagnostic-complete terminus at the framework's structural
reach.  Riemann's RH remains open.  The framework has named the exact live
external interface and the structural reasons that the other surveyed routes
do not advance beyond typed foreclosure or native analog closure.

\section*{Proof Composition}

The proof is by composition of the accepted records:

\begin{itemize}
\item Steps 173--176 and 196 decide Branch C's full-carrier route
numerically through nonzero \(L_{\rho,k}(G)\) values.
\item Steps 177--179, 191--192 identify the shared Burnol \(\kappa\)
blockage for Branches A and B.
\item Steps 180, 183, 187 formalize CTMT, CRCFT, and the Dichotomy.
\item Step 189 derives Bridge Impossibility in partial form.
\item Steps 185, 186, and 190 supply non-CRE native closures.
\item Steps 193--195, 197, and 199 test additional CRE and numerical side
interfaces without finding a Dichotomy refuter.
\item Step 198 integrates the framework audit.
\end{itemize}

These records exhaust the surveyed carrier families and leave only Path 1 or
Path 3 as future Riemann-relevant work.

\section*{Retrospective}

The arc contains 17 major substantive units and 5 synthesis units.  The
synthesis units are CTMT, CRCFT, the Dichotomy, the framework audit, and this
closeout.  The substantive units include projection derivation, numerical
experiments, carrier pivots, theorem derivations, native closure audits, and
external-source audits.

\section*{Next Directions}

\begin{enumerate}
\item External classical research on Burnol \(\kappa\).
\item Search for an in-scope Dichotomy coverage refuter.
\item Integrate CTMT, CRCFT, Dichotomy, and Bridge Impossibility into the
foundational corpus.
\item Test CTMT/CRCFT on other Clay-track carriers.
\item Extend Beurling--Nyman and Branch C numerical sidecars.
\end{enumerate}

\section*{Nonclaim}

This closeout does not prove RH, close \(\Xi_{BC}\), or supply the missing
Burnol projection theorem.  It declares the iteration arc terminus at the
framework's current structural reach.

\end{document}

